\section{Discussion and Conclusion}
\label{sec:conclusion}

Constructive and destructive regimes coexist with learning progress. Around Adam loss spikes, consecutive updates made the network's predictions oscillate along a few low-dimensional recurring functional directions. Constraining large changes in selected hidden representations smoothed training and lowered mean training loss across all three models, although the cap can add additional computational cost. One practical approach would be to apply the cap only when the spike condition in Eq.~\eqref{eq:loss-rise-alignment}, $\rho_t<-\|g_{t-1}\|/\|g_t\|$, warns of a possible loss increase on the next step.

These findings point to a direction for stabilizing training, but the size and scope of our study are limited. We focused on multiple models, modailities and full-batch training with the goal to connect the observed regimes to their theoretical description. We examined how mini-batch sampling affects gradient alignment, but how the representation dynamics and the intervention manifest in large models trained with mini-batches remains open. The cap also leaves the oscillation itself unresolved. It reduced the largest changes in outputs, yet under q25 the Transformer's predictions still reversed direction on nearly nine out of ten updates.

Our net displacement ratios make the inefficiency of gradient-based training explicit. Many updates can produce substantial movement in parameter space or in the learned function, while only a small fraction of that movement remains as net displacement. In the Adam spike windows, predictions retained just 1.7--3.5\% of the distance traveled. Despite the network's many parameters, training appears to explore only a narrow low-dimensional set of functional changes over these windows. For further research, we therefore ask if a hybrid approach that combines the task loss with additional objectives could change the learning dynamics. One would encourage sparsity or compactness in the representations to constrain oscillation. The other would increase the sensitivity of persistently inactive units so they can contribute to learning.
